% ================================================================
% Section: Domain-Specific Applications
% From main (2).tex — Section VII
% ================================================================

\section{Domain-Specific Applications}
\label{sec:applications}

The domains below are where \gls{slm} deployment is most advanced, and each is an
agentic setting in the sense of Section~\ref{sec:agentic}: a workflow of bounded,
repeated language operations---extract a diagnosis code, classify a clause,
complete a function, translate a passage---executed under regulatory or cost
constraints that make a cloud frontier model impractical. Domain fine-tuning of a
small model is the LLM-to-SLM conversion pipeline
(Section~\ref{subsec:agentic_conversion}) applied to a vertical.

\subsection{Healthcare and Clinical NLP}
\label{subsec:applications_healthcare}

Healthcare \gls{nlp} presents unique requirements: high factual accuracy, interpretability, HIPAA/GDPR compliance, and compatibility with clinical workflows. Large general-purpose \glspl{llm} are increasingly restricted from processing identifiable patient data due to data residency regulations, creating a strong pull toward on-premises \glspl{slm}.

The Meerkat family \cite{kim2025meerkat} demonstrates the potential of domain-specific \gls{slm} fine-tuning. Built on a 7B base with instruction fine-tuning on curated clinical QA datasets and chain-of-thought reasoning traces from medical board examinations, Meerkat achieves 85.7\% on USMLE Step~1---surpassing GPT-4 (86.0\% in zero-shot) while being deployable on a single A100 GPU rather than a multi-GPU cloud cluster. Meerkat's clinical reasoning is particularly strong on diagnostic pattern matching, outperforming GPT-4 by 4\% on complex differential diagnosis tasks \cite{kim2025meerkat}.

OptimCLM \cite{hasan2024optimclm} applies the full compression pipeline---knowledge distillation, structured pruning, and 4-bit quantization---to clinical language models, reducing a 7B model to under 1B while preserving 95\% of its accuracy on ICD coding and discharge summary generation tasks. This makes clinical \gls{nlp} feasible on hospital workstations without GPU infrastructure.

The broader \gls{slm} healthcare ecosystem includes BioMedLM (PubMed-trained 2.7B model), clinical note generation models, and radiology-specific models. A key trend is the use of \gls{rag} \cite{lewis2020rag} to anchor \gls{slm} outputs to verified clinical knowledge bases, reducing hallucination rates in patient-facing applications.

\subsection{Legal NLP}
\label{subsec:applications_legal}

Legal document processing---contract review, compliance monitoring, regulatory change analysis, case law search---involves long, precisely structured documents where hallucination is particularly dangerous. \glspl{slm} address legal \gls{nlp} deployment constraints in two ways: (1)~by enabling on-premises deployment to satisfy attorney-client privilege and data security requirements; and (2)~by enabling cost-effective processing of the massive volumes of routine legal work that would be prohibitively expensive with frontier \glspl{llm}.

Pilot deployments of domain-adapted \glspl{slm} (3B--7B range, fine-tuned on contract corpora) have demonstrated 30--40\% reductions in administrative processing time for standard contract review workflows. Models fine-tuned on specific regulatory domains (GDPR, SOC~2, FINRA) consistently outperform general-purpose GPT-4 on domain-specific compliance classification tasks, owing to the depth of in-domain knowledge concentration.

\subsection{Code Generation}
\label{subsec:applications_code}

Code-focused \glspl{slm} have matured into a category distinct from general \glspl{slm}, with specialized pre-training corpora (GitHub, Stack Overflow, documentation), code-aware tokenization, and evaluation on function-level benchmarks (HumanEval, MBPP, SWE-Bench).

\begin{description}
    \item[Code LLaMA \cite{roziere2023codellama}] Meta's family of code models fine-tuned on LLaMA-2 (7B, 13B, 34B). Code LLaMA 7B achieves 33.5\% HumanEval pass@1.

    \item[StarCoder2 \cite{lozhkov2024starcoder2}] BigCode's 3B/7B/15B models trained on The Stack v2 (600+ programming languages). StarCoder2-7B achieves 46.3\% HumanEval pass@1.

    \item[DeepSeek-Coder \cite{zhu2024deepseek_coder}] A 1B--33B family trained from scratch on 2T code tokens. The 6.7B variant achieves 47.6\% HumanEval pass@1, outperforming GPT-3.5 on code completion.

    \item[Codestral \cite{mistral2024codestral}] Mistral's 22B code model trained on 80+ programming languages with fill-in-the-middle (FIM) training for IDE integration.
\end{description}

\subsection{Education and Multilingual NLP}
\label{subsec:applications_multilingual}

Multilingual \glspl{slm} address the linguistic coverage gap in AI assistants: the majority of \gls{llm} pre-training data is English, leaving speakers of lower-resource languages with significantly degraded performance. Aya~8B \cite{ustun2024aya} is fine-tuned on 204 million multilingual instructions spanning 101 languages, with strong performance across African, South Asian, and Southeast Asian languages that are severely underrepresented in standard benchmarks.

Qwen2.5~7B \cite{qwen2024qwen25} is pre-trained on 18T tokens with strong multilingual coverage including Chinese, Japanese, Korean, Arabic, and multiple European languages, serving as the de facto open-weight \gls{slm} for multilingual enterprise deployments. Gemma~3's multilingual extension \cite{team2025gemma3} adds instruction-tuning data for 40+ languages at the 4B and 12B scales.

Educational applications include on-device tutoring systems (leveraging Phi-3 or SmolLM on tablets), interactive language learning assistants with real-time feedback, and low-connectivity educational tools for rural and developing-world deployments. The key requirement for educational \glspl{slm} is not just multilingual capability but also low hallucination rates and age-appropriate safety filtering.
